\subsection{The sixth conductor is impossible}

We combine the integral canonical class with the same slope bound
to exclude conductor six, including the singular characteristic-two case.

\begin{proposition}\label{prop:conductor-six-excluded}
Under the hypotheses of Theorem~\ref{thm:conductor-equality-structure},
the conductor $a=6$ is impossible in every characteristic.
\end{proposition}

\begin{proof}
Suppose $a=6$, and retain $d=Y^2$ and $\epsilon\in\{0,1\}$
from the structure theorem. On the minimal smooth normalization
resolution $T$, the Cartier canonical identity and Chern numbers are
\[
 K_T=5L,\qquad K_T^2=25pd,\qquad
 c_2(T)=-6pd+2\epsilon,\qquad
 12\chi(\OO_T)=19pd+2\epsilon.
\]
The surface $T$ is minimal of general type, its second Chern number
is negative, and its Euler characteristic is positive.
The negative-Chern-number theorem supplies an Albanese fibration
of arithmetic genus $g\geq2$
\cite[Theorems~2 and~7, pp.~264--266]{ShepherdBarronGeography}.
Adjunction on a fibre gives $5\mid2g-2$, so $g\equiv1\pmod5$
and $g\geq6$. This uses arithmetic genus and does not require a
smooth general fibre.

For $g\geq6$, Gu--Sun--Zhou's nonhyperelliptic bound $n(g)$
dominates the hyperelliptic bound $h(g)$ and decreases with $g$.
Indeed,
\[
 n(g)-h(g)=
 \frac{4(g-1)(3g+2)(g^2-4g+2)}
 {g(g^2-4g+5)(g^2-g-1)}>0,
\]
\[
 n'(g)=-\frac{4\{5g^2((g-3)^2-2)+4(g-1)(3g-5)\}}
 {(g^3-4g^2+5g)^2}<0.
\]
Thus their theorem gives
\[
 \frac{K_T^2}{\chi(\OO_T)}\leq n(6)=\frac{800}{51}
 \qquad\text{\cite[Theorem~3.2, pp.~622--623]{GuSunZhou}.}
\]

If $\epsilon=0$, the ratio is $300/19$, and
$300/19-800/51=100/969>0$, a contradiction.
If $\epsilon=1$, then $p=2$, and the ratio is
$300d/(19d+1)$. The same inequality gives
\[
 15300d\leq15200d+800,
\]
so $d\leq8$.
On the original smooth surface, $K_X=11Y$ and
$c_2(X)=-12d+1$. The latter follows from the structure theorem's
quotient calculation: the blowup of the single internal center
increases $c_2$ by one and has the same $c_2$ as $T$.
Noether's formula on $X$ therefore yields
\[
 12\chi(\OO_X)=109d+1.
\]
Hence $d\equiv11\pmod{12}$, so $d\geq11$, contradicting
$d\leq8$. This excludes the internal case as well.
\end{proof}
